% =============================================================
%  AI Surrogate Modeling for Accelerating Chaotic Lorenz
%  Simulation using OpenMP and CUDA
%  Author : Jovan Moris  |  CS23B2058
%  Date   : March 2026
% =============================================================

\documentclass[11pt,a4paper]{article}

% - Fonts -
\usepackage[T1]{fontenc}
\usepackage[utf8]{inputenc}
\usepackage{palatino}
\usepackage{mathpazo}
\usepackage{microtype}

% - Mathematics -
\usepackage{amsmath,amssymb}

% - Page geometry -
\usepackage[top=25mm,bottom=25mm,left=28mm,right=28mm,
            headheight=14pt]{geometry}

% - Headers & footers -
\usepackage{fancyhdr}
\pagestyle{fancy}
\fancyhf{}
\renewcommand{\headrulewidth}{0.4pt}
\fancyhead[L]{\small\textit{AI Surrogate Modeling --- Lorenz System}}
\fancyhead[R]{\small\textit{Jovan Moris D - CS23B2058}}
\fancyfoot[C]{\thepage}

% - Colours -
\usepackage[dvipsnames,table]{xcolor}
\definecolor{accent}{RGB}{0,76,153}
\definecolor{lightblue}{RGB}{235,242,255}
\definecolor{darkgray}{RGB}{80,80,80}

% - Graphics -
\usepackage{graphicx}
\usepackage{float}
\usepackage{subcaption}
\usepackage[labelfont={bf,color=accent},font=small,
            justification=centering]{caption}

% - Tables -
\usepackage{booktabs}
\usepackage{multirow}
\usepackage{array}
\usepackage{tabularx}
\newcolumntype{C}{>{\centering\arraybackslash}X}

% - Lists -
\usepackage{enumitem}
\setlist[itemize]{noitemsep,topsep=3pt,leftmargin=1.4em}
\setlist[enumerate]{noitemsep,topsep=3pt,leftmargin=1.6em}

% - Code listings -
\usepackage{listings}
\lstdefinestyle{mycode}{
  basicstyle=\ttfamily\footnotesize,
  keywordstyle=\color{accent}\bfseries,
  commentstyle=\color{OliveGreen}\itshape,
  stringstyle=\color{BrickRed},
  numbers=left, numberstyle=\tiny\color{gray},
  frame=single, rulecolor=\color{gray!40},
  breaklines=true, captionpos=b,
  backgroundcolor=\color{gray!6},
  xleftmargin=1.6em, xrightmargin=0.4em,
  aboveskip=6pt, belowskip=4pt,
  showstringspaces=false
}
\lstset{style=mycode, language=Python}

% - Section headings -
\usepackage{titlesec}
\titleformat{\section}
  {\large\bfseries\color{accent}}{\thesection}{1em}{}[\color{accent}\titlerule]
\titleformat{\subsection}
  {\normalsize\bfseries\color{accent!85!black}}{\thesubsection}{1em}{}
\titleformat{\subsubsection}
  {\normalsize\bfseries\color{darkgray}}{\thesubsubsection}{1em}{}
\titlespacing{\section}{0pt}{14pt}{6pt}
\titlespacing{\subsection}{0pt}{10pt}{4pt}

% - Spacing -
\setlength{\parskip}{5pt}
\setlength{\parindent}{0pt}
\renewcommand{\baselinestretch}{1.13}

% - Hyperlinks -
\usepackage[colorlinks=true,linkcolor=accent,
            citecolor=accent,urlcolor=accent]{hyperref}

% - Bibliography -
\usepackage[numbers,sort&compress]{natbib}

% - Convenience commands -
\newcommand{\good}[1]{{\color{OliveGreen}\textbf{#1}}}
\newcommand{\bad}[1]{{\color{BrickRed}\textit{#1}}}

% =============================================================
\begin{document}

% - Title page -
\begin{titlepage}
\centering
\vspace*{1.5cm}
{\color{accent}\rule{\linewidth}{2pt}}\\[0.4cm]
{\LARGE\bfseries\color{accent}
AI Surrogate Modeling for Accelerating\\[0.3em]
Chaotic Lorenz Simulation\\[0.3em]
using OpenMP and CUDA}\\[0.4cm]
{\color{accent}\rule{\linewidth}{1pt}}\\[1.2cm]
{\large AI-HPC Project}\\[1.8cm]

\begin{tabular}{rl}
  \textbf{Author}      & Jovan Moris D\\[4pt]
  \textbf{Roll No.}    & CS23B2058 \\[4pt]
  \textbf{Department}  & Computer Science \& Engineering \\[4pt]
  \textbf{Course}      & High Performance Computing \\[4pt]
  \textbf{Date}        & March 2026 \\
\end{tabular}

\vfill

\begin{abstract}
\noindent
This report presents a complete implementation of a hybrid HPC and AI pipeline
for accelerating the simulation of the Lorenz chaotic dynamical system. The
system architecture works in two complementary stages: \textbf{(1)} parallel
computation using OpenMP (multi-core CPU) and CUDA (GPU) to rapidly generate
trajectories up to a defined checkpoint, and \textbf{(2)} trained AI surrogate
models that take over from the checkpoint to predict the remaining trajectory
without further numerical integration. This design correctly reflects the
complementary roles of HPC and AI --- HPC accelerates the physics-based simulation
while AI replaces it beyond the predictability horizon.

The HPC component delivers the project's primary contribution: OpenMP achieves
\good{1682$\times$} speedup and CUDA achieves \good{3299$\times$} speedup over
the Python RK4 baseline. For the AI surrogate component, four models are
implemented and evaluated: MLP, improved LSTM, Temporal Convolutional Network
(TCN), and Echo State Network (ESN). TCN achieves a single-step RMSE of
\good{0.0128} and ESN achieves \good{0.0134} --- representing a \good{97.6\%}
improvement over the original baseline RMSE of 0.5256. The ESN is particularly
notable, achieving near-identical accuracy to the TCN while training in just
\good{33 seconds} compared to TCN's 159 minutes.
\end{abstract}

\vspace{0.5cm}
{\color{accent}\rule{\linewidth}{2pt}}
\end{titlepage}

% - TOC -
\tableofcontents
\newpage
\listoffigures
\listoftables
\newpage

% =============================================================
\section{Introduction}
% =============================================================

\subsection{Motivation and Problem Statement}

The Lorenz system, introduced by Edward Lorenz in 1963~\cite{lorenz1963}, is governed
by three coupled ordinary differential equations:

\begin{align}
  \frac{dx}{dt} &= \sigma(y - x) \label{eq:lx}\\
  \frac{dy}{dt} &= x(\rho - z) - y \label{eq:ly}\\
  \frac{dz}{dt} &= xy - \beta z \label{eq:lz}
\end{align}

With standard chaotic parameters $\sigma=10$, $\rho=28$, $\beta=\tfrac{8}{3}$,
trajectories never repeat and are exquisitely sensitive to initial conditions ---
the famous \textit{butterfly effect}. Simulating large ensembles of such trajectories
is computationally demanding, making this system an ideal testbed for HPC and AI
acceleration.

\subsection{System Architecture: The Correct HPC + AI Design}

A key insight of this project is that HPC and AI serve \textbf{complementary, not
competing, roles}:

\begin{enumerate}
  \item \textbf{Phase 1 --- HPC Simulation to Checkpoint:} OpenMP and CUDA rapidly
        generate trajectories up to a defined checkpoint using parallelised RK4
        integration. This phase exploits parallel structure: each
        trajectory is independent, requiring no inter-thread communication.
  \item \textbf{Phase 2 --- AI Surrogate from Checkpoint:} A trained surrogate model
        (MLP or LSTM or TCN or ESN) takes the checkpoint state as input and predicts the remaining
        trajectory without further numerical integration. This replaces the expensive
        per-step ODE solving with a single fast forward pass.
\end{enumerate}

This design is both physically motivated and computationally sound. The HPC component
handles the deterministic early phase accurately; the AI component handles the longer
horizon where the Lyapunov exponent ($\lambda \approx 0.906$) limits precise
predictability anyway.


\subsection{Problem Statement}

\subsubsection{Why Lorenz Simulation is Expensive}

Generating a single Lorenz trajectory requires solving three coupled ODEs at every
time step. Using RK4 with $\Delta t = 0.01$ over $T = 25$ seconds requires 2500
sequential integration steps per trajectory. For a single trajectory this is
manageable, but real-world applications require ensembles of thousands to millions
of trajectories:

\begin{itemize}
  \item \textbf{Ensemble forecasting:} Weather models require $10^3$--$10^6$
        trajectories to quantify forecast uncertainty.
  \item \textbf{Uncertainty quantification:} Monte Carlo methods for sensitivity
        analysis require thousands of independent runs.
  \item \textbf{Parameter sweeps:} Bifurcation analysis sweeps $\sigma$, $\rho$,
        $\beta$ across a grid, requiring a full simulation per grid point.
  \item \textbf{Optimal control:} Real-time trajectory optimisation requires
        repeated fast simulation to evaluate candidate control actions.
  \item \textbf{Data-driven surrogate training:} Training neural surrogates requires
        generating large datasets of trajectories as ground truth.
\end{itemize}

In pure Python, generating 100 trajectories takes $\approx 3.4$ seconds. Scaling to
$10^5$ trajectories would require $\approx 57$ minutes --- far too slow for real-time
or interactive applications. This motivates both HPC acceleration and AI surrogates.

\subsubsection{Why Standard AI Surrogates Struggle with Chaos}

Chaotic systems pose a unique challenge for data-driven models. The positive Lyapunov
exponent ($\lambda \approx 0.906$) means prediction errors grow exponentially as
$\varepsilon(t) = \varepsilon_0 e^{\lambda t}$. This imposes a hard
\textbf{predictability horizon} of $\approx 1/\lambda \approx 1.1$ time units
($= 110$ steps at $\Delta t=0.01$) beyond which \textit{any} deterministic model
--- no matter how accurate --- will diverge from the true trajectory.

This is not a failure of the model; it is a physical law. The correct response is not
to expect long-horizon pointwise accuracy, but to instead ask:
\begin{enumerate}
  \item Can the surrogate predict accurately within the Lyapunov horizon?
  \item Does the surrogate remain on the correct attractor after divergence?
  \item Can the surrogate replace expensive simulation for rapid approximate screening?
\end{enumerate}

\subsubsection{The Hybrid HPC + AI Solution}

This project proposes a two-phase pipeline that correctly addresses both challenges:

\begin{enumerate}
  \item \textbf{Phase 1 (HPC):} Generate trajectories rapidly and accurately up to
        a checkpoint using OpenMP or CUDA parallelised RK4. Since trajectories are
        independent, this is embarrassingly parallel --- perfect for multi-core and
        GPU architectures.
  \item \textbf{Phase 2 (AI):} From the checkpoint state, a trained surrogate model
        (TCN or ESN) predicts the remaining trajectory in a single forward pass,
        without further numerical integration.
\end{enumerate}

The speedup comes from two directions: HPC reduces the cost of accurate simulation,
and AI eliminates it entirely beyond the checkpoint. Together they enable trajectory
generation at speeds impractical with either approach alone.

\subsection{Research Objectives}

\begin{enumerate}
  \item Implement OpenMP and CUDA parallelisation for RK4-based trajectory generation
        and benchmark speedups.
  \item Train and evaluate four AI surrogate models: MLP, improved LSTM, TCN, and ESN.
  \item Demonstrate the complementary HPC+AI pipeline on the Lorenz system.
  \item Honestly analyse both achievements and limitations.
\end{enumerate}


% =============================================================
\section{Mathematical Foundation}
% =============================================================

\subsection{Lorenz System Properties}

The system \eqref{eq:lx}--\eqref{eq:lz} with standard parameters has:
\begin{itemize}
  \item Largest Lyapunov exponent $\lambda \approx 0.906$ --- trajectories diverge
        as $\varepsilon e^{\lambda t}$.
  \item Predictability horizon $\approx 1/\lambda \approx 1.1$ time units
        ($= 110$ steps at $\Delta t=0.01$).
  \item Strange attractor with Hausdorff dimension $\approx 2.06$.
\end{itemize}

\subsection{Fourth-Order Runge-Kutta Integration}

\begin{align}
  \mathbf{k}_1 &= f(\mathbf{y}_n), \quad
  \mathbf{k}_2 = f\!\left(\mathbf{y}_n + \tfrac{\Delta t}{2}\mathbf{k}_1\right)\\
  \mathbf{k}_3 &= f\!\left(\mathbf{y}_n + \tfrac{\Delta t}{2}\mathbf{k}_2\right),
  \quad \mathbf{k}_4 = f(\mathbf{y}_n + \Delta t\,\mathbf{k}_3)\\
  \mathbf{y}_{n+1} &= \mathbf{y}_n +
      \frac{\Delta t}{6}(\mathbf{k}_1 + 2\mathbf{k}_2 + 2\mathbf{k}_3 + \mathbf{k}_4)
\end{align}

\subsection{Supervised Learning Formulation}

Surrogate modelling is cast as sliding-window regression:
\begin{equation}
  \mathcal{D} = \bigl\{(\mathbf{X}^{(i)},\,\mathbf{y}^{(i)})\bigr\}_{i=1}^{N},
  \quad \mathbf{X}^{(i)} \in \mathbb{R}^{L \times 3},\quad
  \mathbf{y}^{(i)} \in \mathbb{R}^{3}
\end{equation}

with $L=50$ (lookback window), z-score normalisation applied to training split only.

% =============================================================
\section{Implementation}
% =============================================================

\subsection{OpenMP C++ Parallelisation}

The C++ solver parallelises over trajectories using \texttt{\#pragma omp parallel for}.
Each trajectory is fully independent --- no shared state, no synchronisation needed.

\begin{lstlisting}[language=C++,
  caption={OpenMP parallel trajectory loop --- \texttt{src/openmp\_lorenz.cpp}}]
#pragma omp parallel for schedule(dynamic)
for (int traj = 0; traj < num_trajectories; traj++) {
    double x = init_x[traj], y = init_y[traj], z = init_z[traj];
    for (int step = 0; step < total_steps; step++) {
        rk4_step(&x, &y, &z, dt);
        trajectories[traj][step][0] = x;
        trajectories[traj][step][1] = y;
        trajectories[traj][step][2] = z;
    }
}
\end{lstlisting}

Compilation: \texttt{g++ -O3 -fopenmp -march=native src/openmp\_lorenz.cpp}.

\newpage

\subsection{CUDA GPU Implementation}

The CUDA kernel assigns one thread per trajectory. Each thread runs the complete RK4
loop independently, operating entirely in registers with no global memory
synchronisation required.

\begin{lstlisting}[language=C,
  caption={CUDA kernel --- \texttt{cuda/cuda\_lorenz.cu}}]
__global__ void lorenzKernel(double* trajectories,
                              int total_steps, double dt,
                              int num_trajectories) {
    int tid = blockIdx.x * blockDim.x + threadIdx.x;
    if (tid >= num_trajectories) return;

    double x = init_x[tid], y = init_y[tid], z = init_z[tid];
    for (int step = 0; step < total_steps; step++) {
        rk4_step_cuda(&x, &y, &z, dt);
        int idx = tid * total_steps * 3 + step * 3;
        trajectories[idx]   = x;
        trajectories[idx+1] = y;
        trajectories[idx+2] = z;
    }
}
\end{lstlisting}

Compilation: \texttt{nvcc -O3 cuda/cuda\_lorenz.cu}.

\subsection{AI Surrogate Models}

\subsubsection{MLP Surrogate}
Flattened input ($50 \times 3 = 150$ features), architecture $150 \to 128 \to 64
\to 32 \to 3$. Acts as a flat baseline with no temporal inductive bias.

\subsubsection{Improved LSTM}
Hidden size 256, 3 layers, dropout 0.2. Processes window sequentially; learns
temporal dependencies unavailable to MLP.

\subsubsection{Temporal Convolutional Network (TCN)}
Four dilated causal convolution blocks with dilation factors 1, 2, 4, 8. Each block:
CausalConv1d $\to$ BatchNorm $\to$ ReLU $\to$ Dropout $\to$ residual connection.
Receptive field grows exponentially without vanishing gradients.

\newpage 


\subsubsection{Echo State Network (ESN)}
A fixed random reservoir acts as a nonlinear dynamical memory. Only the linear
readout weights are trained via ridge regression --- no backpropagation, no epochs.

\begin{lstlisting}[caption={ESN reservoir update and readout training}]
# Reservoir update (leaky integrator)
x = (1-lr)*x + lr*np.tanh(W @ x + W_in @ u)

# Readout training - ridge regression only
readout = Ridge(alpha=1e-4)
readout.fit(reservoir_states[washout:], data[washout+1:])
\end{lstlisting}

Key hyperparameters: \texttt{n\_reservoir=500}, \texttt{spectral\_radius=0.95}
(must be $<1.0$ for stability), \texttt{leaking\_rate=0.5}, \texttt{washout=200}.

\subsection{Model Architecture Diagrams}

Figures \ref{fig:mlp_arch}--\ref{fig:esn_arch} illustrate the architecture of each
surrogate model. These diagrams show the data flow, layer dimensions, and key
design choices for each approach.

\begin{figure}[H]
  \centering
  \includegraphics[width=\textwidth]{arch_mlp.png}
  \caption{MLP surrogate architecture. The input window (50 timesteps x 3 dimensions
           = 150 features) is flattened and passed through three fully connected
           hidden layers with ReLU activation and dropout. No temporal structure
           is exploited.}
  \label{fig:mlp_arch}
\end{figure}

\begin{figure}[H]
  \centering
  \includegraphics[width=\textwidth]{arch_lstm.png}
  \caption{Improved LSTM surrogate architecture. Three stacked LSTM layers
           (hidden=256) process the input sequence step by step, maintaining
           hidden state $ and cell state $ across timesteps. Only the last
           hidden state feeds the output head.}
  \label{fig:lstm_arch}
\end{figure}

\begin{figure}[H]
  \centering
  \includegraphics[width=\textwidth]{arch_tcn.png}
  \caption{TCN surrogate architecture. Four dilated causal convolution blocks
           with dilation factors 1, 2, 4, 8 achieve a receptive field of 31+
           timesteps without any recurrence. Residual connections prevent
           vanishing gradients.}
  \label{fig:tcn_arch}
\end{figure}

\begin{figure}[H]
  \centering
  \includegraphics[width=\textwidth]{arch_esn.png}
  \caption{ESN (Echo State Network) architecture. The reservoir (500 neurons,
           spectral radius=0.95) is fixed and random --- never trained. Only the
           linear readout layer is fitted via ridge regression. This makes ESN
           training orders of magnitude faster than gradient-based methods.}
  \label{fig:esn_arch}
\end{figure}

\subsection{Training Configuration}

\begin{center}
\rowcolors{2}{lightblue}{white}
\begin{tabular}{lll}
  \toprule
  \rowcolor{accent!15}
  \textbf{Parameter} & \textbf{Deep Models} & \textbf{ESN} \\
  \midrule
  Trajectories    & 500             & 500 (30\% used) \\
  Lookback $L$    & 50 steps        & Full sequence \\
  Optimiser       & Adam            & Ridge regression \\
  Loss            & MSE             & Least squares \\
  Epochs          & 50/12/80 (M/L/T)& N/A \\
  Grad clipping   & max\_norm=1.0   & N/A \\
  LR scheduler    & ReduceLROnPlateau & N/A \\
  \bottomrule
\end{tabular}
\end{center}

\newpage

% =============================================================
\section{Experimental Setup}
% =============================================================

\begin{table}[H]
\centering
\caption{Simulation parameters used throughout all experiments.}
\label{tab:params}
\rowcolors{2}{lightblue}{white}
\begin{tabular}{lll}
  \toprule
  \rowcolor{accent!15}
  \textbf{Parameter} & \textbf{Value} & \textbf{Description} \\
  \midrule
  $\sigma$, $\rho$, $\beta$  & 10, 28, 8/3  & Standard chaotic parameters \\
  Trajectories (HPC)         & 100           & For speedup benchmarking \\
  Trajectories (AI)          & 500           & For surrogate training \\
  Simulation time $T$        & 25.0 s        & Total integration time \\
  Time step $\Delta t$       & 0.01 s        & RK4 step size \\
  Steps per trajectory       & 2500          & $T / \Delta t$ \\
  Train / Val / Test         & 70\% / 15\% / 15\% & Dataset split \\
  Random seed                & 42            & Reproducibility \\
  HPC platform               & Local Linux   & \texttt{g++}, \texttt{nvcc} \\
  AI platform                & Kaggle T4 GPU & PyTorch 2.9, CUDA 12.6 \\
  \bottomrule
\end{tabular}
\end{table}

% =============================================================
\section{Results}
% =============================================================

\subsection{Lorenz Attractor Visualisation}

\begin{figure}[H]
  \centering
  \includegraphics[width=0.72\textwidth]{01_lorenz_attractor_3d.png}
  \caption{3D Lorenz strange attractor computed via RK4 integration
           ($\sigma=10$, $\rho=28$, $\beta=8/3$, $\Delta t=0.01$, $T=25$\,s).
           The characteristic double-lobe butterfly shape confirms correct
           chaotic regime behaviour.}
  \label{fig:attractor}
\end{figure}

\begin{figure}[H]
  \centering
  \includegraphics[width=\textwidth]{05_phase_space_projections.png}
  \caption{2D phase-space projections of the Lorenz attractor:
           $X$--$Y$ (left), $X$--$Z$ (centre), $Y$--$Z$ (right).
           Green dot: trajectory start; red cross: trajectory end.}
  \label{fig:phase}
\end{figure}

\begin{figure}[H]
  \centering
  \includegraphics[width=\textwidth]{06_trajectory_evolution.png}
  \caption{Time-series evolution of $X$, $Y$, $Z$ over 500 integration steps.
           The irregular, non-repeating oscillations confirm chaotic dynamics.}
  \label{fig:stateevol}
\end{figure}

\subsection{HPC Performance: Phase 1 (Simulation to Checkpoint)}

The primary purpose of OpenMP and CUDA in this pipeline is to rapidly generate
trajectories up to the checkpoint using parallelised RK4 integration. Since each
trajectory is independent, this is an embarrassingly parallel workload that maps
perfectly to both multi-core and many-core architectures.

\begin{table}[H]
\centering
\caption{HPC benchmarking: 100 trajectories, 2500 steps each.
         Speedup relative to Python RK4 baseline.}
\label{tab:hpc}
\rowcolors{2}{lightblue}{white}
\begin{tabular}{lrrr}
  \toprule
  \rowcolor{accent!15}
  \textbf{Method} & \textbf{Time (ms)} &
  \textbf{Speedup} & \textbf{Throughput (traj/s)} \\
  \midrule
  Python RK4 (baseline) & 3365.23 & 1.00$\times$              & 29.71    \\
  Python Euler          & 762.45  & 4.41$\times$               & 131.15   \\
  OpenMP (C++)          & 2.00    & \good{1682.6$\times$}      & 50{,}000 \\
  CUDA (GPU)            & 1.02    & \good{3299.2$\times$}      & 98{,}039 \\
  \bottomrule
\end{tabular}
\end{table}

\begin{figure}[H]
  \centering
  \includegraphics[width=\textwidth]{04_performance_comparison.png}
  \caption{Left: absolute execution times for all methods. Right: speedup factors
           relative to Python RK4. OpenMP and CUDA deliver the dominant speedups
           for Phase 1 trajectory generation up to the checkpoint.}
  \label{fig:hpcperf}
\end{figure}

The CUDA/OpenMP speedup ratio ($\approx 2\times$) at $N=100$ trajectories reflects
the relatively small problem size for GPU utilisation. At $N=10{,}000$ trajectories,
this ratio would widen to $10$--$50\times$ as GPU parallelism is better exploited.

\subsection{AI Surrogate Results: Phase 2 (Checkpoint to Future)}

Once HPC delivers the checkpoint state, the AI surrogate takes over. Four models
were trained and evaluated on 500 trajectories with lookback window $L=50$.

\subsubsection{Training Convergence}

\begin{figure}[H]
  \centering
  \includegraphics[width=\textwidth]{1773840048444_02_training_history.png}
  \caption{Training and validation MSE loss (log scale) for MLP (left),
           improved LSTM (centre), and TCN (right) over their respective
           epochs. TCN shows the fastest and deepest convergence.}
  \label{fig:training}
\end{figure}

\subsubsection{Single-Step Prediction Accuracy}

\begin{table}[H]
\centering
\caption{Single-step prediction accuracy on the test set (normalised units).
         Green cells indicate best value in each column.}
\label{tab:singlestep}
\rowcolors{2}{lightblue}{white}
\begin{tabular}{lrrrr}
  \toprule
  \rowcolor{accent!15}
  \textbf{Model} & \textbf{RMSE} & \textbf{MAE} &
  \textbf{Improvement vs Baseline} & \textbf{Train Time} \\
  \midrule
  MLP (original baseline) & 0.5256 & 0.4574 & ---              & 200 epochs \\
  MLP (improved)          & 0.4027 & 0.3184 & +23.4\%           & 33 min \\
  LSTM (improved)         & 0.2179 & 0.1643 & +58.5\%           & 41 min \\
  TCN                     & \good{0.0128} & \good{0.0079} & \good{+97.6\%} & 159 min \\
  ESN                     & 0.0134 & 0.0096 & +97.5\%           & \good{33 sec} \\
  \bottomrule
\end{tabular}
\end{table}

\begin{figure}[H]
  \centering
  \includegraphics[width=\textwidth]{07_rmse_mae_comparison.png}
  \caption{Single-step RMSE (left) and MAE (right) for all four models.
           TCN and ESN achieve near-identical accuracy, both reducing error
           by $\approx$97.5\% over the original baseline.}
  \label{fig:rmse}
\end{figure}

\begin{figure}[H]
  \centering
  \includegraphics[width=0.85\textwidth]{10_results_table.png}
  \caption{Complete results summary table. Green cells indicate best
           performance in each metric category.}
  \label{fig:table}
\end{figure}

\subsubsection{Training Time vs Accuracy Tradeoff}

\begin{figure}[H]
  \centering
  \includegraphics[width=0.85\textwidth]{08_time_vs_accuracy.png}
  \caption{Training time vs RMSE tradeoff. The ESN achieves essentially
           the same accuracy as TCN (0.0134 vs 0.0128) while training in
           33 seconds rather than 159 minutes --- a 290$\times$ faster
           training time for equivalent accuracy.}
  \label{fig:tradeoff}
\end{figure}

\subsubsection{Improvement Over Original Baseline}

\begin{figure}[H]
  \centering
  \includegraphics[width=0.85\textwidth]{11_improvement_over_baseline.png}
  \caption{Percentage RMSE improvement of each model over the original
           baseline (RMSE = 0.5256). TCN and ESN both achieve $>$97\%
           improvement.}
  \label{fig:improvement}
\end{figure}

\subsubsection{Multi-Step Rollout Analysis}

The most demanding test is autoregressive rollout: each model's prediction is fed
back as input for subsequent steps. This tests whether the surrogate can sustain
trajectory quality beyond single-step prediction.

\begin{table}[H]
\centering
\caption{200-step autoregressive rollout L2 error (original units).}
\label{tab:rollout}
\rowcolors{2}{lightblue}{white}
\begin{tabular}{lrr}
  \toprule
  \rowcolor{accent!15}
  \textbf{Model} & \textbf{Mean L2} & \textbf{Max L2} \\
  \midrule
  MLP             & 9.5357  & 28.1619 \\
  LSTM (improved) & 8.5379  & 22.1789 \\
  ESN             & 5.8002  & 26.0734 \\
  TCN             & \good{2.6097}  & \good{13.7378} \\
  \bottomrule
\end{tabular}
\end{table}

\begin{figure}[H]
  \centering
  \includegraphics[width=\textwidth]{09_rollout_error_comparison.png}
  \caption{200-step rollout mean L2 (solid) and max L2 (hatched) for all
           models. TCN achieves the lowest mean error (2.61), nearly
           half that of ESN (5.80) and one-third of MLP/LSTM.}
  \label{fig:rollouterr}
\end{figure}

\begin{figure}[H]
  \centering
  \includegraphics[width=\textwidth]{1773840048444_04_error_horizon.png}
  \caption{L2 error vs prediction step (linear and log scales) for all
           models over 200-step rollout. TCN maintains the lowest error
           profile; ESN shows stable bounded growth.}
  \label{fig:horizon}
\end{figure}

\begin{figure}[H]
  \centering
  \includegraphics[width=\textwidth]{1773840048444_05_rollout_comparison.png}
  \caption{Multi-step rollout trajectory comparison (100 steps).
           Black: RK4 reference. All models track closely for the first
           $\approx$40 steps; TCN and ESN maintain trajectory shape
           longest before diverging due to the Lyapunov instability.}
  \label{fig:rollout}
\end{figure}

\subsubsection{ESN Attractor Stability}

\begin{figure}[H]
  \centering
  \includegraphics[width=\textwidth]{1773840048445_06_phase_space_esn.png}
  \caption{3D phase space comparison over 1000 steps: RK4 reference (left)
           vs ESN rollout (right). The ESN trajectory remains confined to
           the correct attractor region, confirming that the reservoir has
           learned the global attractor geometry even when individual
           trajectory details diverge.}
  \label{fig:esnattractor}
\end{figure}

\begin{figure}[H]
  \centering
  \includegraphics[width=\textwidth]{12_phase_space_2d.png}
  \caption{2D phase-space projections of the Lorenz attractor used for
           surrogate training and evaluation.}
  \label{fig:phase2d}
\end{figure}

% =============================================================
\section{Discussion}
% =============================================================

\subsection{HPC Results --- Primary Contribution}

OpenMP (\good{1682$\times$}) and CUDA (\good{3299$\times$}) deliver the project's
definitive achievement. Both are real, reproducible, and correctly motivated: they
accelerate the physics-based simulation stage that must be accurate. The
embarrassingly parallel structure of multi-trajectory generation maps perfectly to
both architectures.

The CUDA/OpenMP ratio of $\approx$2$\times$ at $N=100$ is expected --- GPU overhead
dominates at small problem sizes. For $N=10^4$, the GPU advantage would be $10$--$50\times$.

\subsection{AI Surrogate Results}

\subsubsection{TCN: Best Accuracy}
TCN achieves RMSE 0.0128 --- a 97.6\% improvement over the baseline. Dilated causal
convolutions capture the temporal structure of the Lorenz attractor across the full
$L=50$ lookback window without vanishing gradient issues. Its rollout mean L2 of 2.61
is the lowest of all models, confirming it as the best surrogate for sustained
prediction.

\subsubsection{ESN: Best Efficiency}
ESN achieves RMSE 0.0134 --- essentially identical accuracy to TCN --- while training
in 33 seconds. This is 290$\times$ faster training than TCN for comparable accuracy.
The ESN's reservoir, being itself a nonlinear dynamical system, naturally resonates
with the Lorenz attractor's structure. Figure~\ref{fig:esnattractor} confirms the ESN
stays on the correct attractor even after 1000 rollout steps.

\subsubsection{Why MLP and LSTM Underperform}
MLP (RMSE 0.4027) lacks temporal inductive bias --- it treats the window as flat
features with no sequential ordering. LSTM (RMSE 0.2179) is limited by the small
epoch budget (12 epochs due to $\approx$3.4 min/epoch on T4) and the difficulty of
backpropagating through long sequences in chaotic systems.

\subsubsection{Rollout Limitations and Physical Explanation}
All models diverge in long-horizon rollout. This is not a modelling failure --- it is
a physical consequence of the positive Lyapunov exponent ($\lambda \approx 0.906$).
Prediction errors grow as $\varepsilon e^{0.906 t}$, limiting reliable prediction to
$\approx 110$ steps regardless of model quality. The ESN's bounded rollout growth
(Figure~\ref{fig:horizon}) and correct attractor confinement (Figure~\ref{fig:esnattractor})
suggest it has captured the global attractor geometry better than other models.

\subsection{Sensitivity Analysis}

\textbf{Lookback window} ($L=10 \to 50$): Increasing from 10 to 50 steps captures
$\approx$45\% of the Lyapunov time, significantly improving LSTM and enabling TCN to
leverage its dilated receptive field.

\textbf{Trajectories} (100 $\to$ 500): Better attractor coverage improves
generalisation. With the CUDA solver, generating 500 trajectories takes
$<$1 second --- making large datasets essentially free.

\textbf{ESN spectral radius}: Must be $<1.0$ for echo state property (stability).
The value 0.95 provides strong memory while maintaining stability.

% =============================================================
\section{Conclusions}
% =============================================================

This project implements and evaluates a complete hybrid HPC+AI pipeline for the
Lorenz system, with the following key conclusions:

\begin{enumerate}
  \item \textbf{HPC is the primary contribution.} OpenMP (\good{1682$\times$}) and
        CUDA (\good{3299$\times$}) deliver exceptional speedups for Phase 1
        trajectory generation. The embarrassingly parallel structure maps perfectly
        to both architectures.

  \item \textbf{TCN is the best surrogate.} RMSE 0.0128, mean rollout L2 of 2.61,
        and the best long-horizon trajectory tracking. Dilated causal convolutions
        are well-matched to chaotic temporal structure.

  \item \textbf{ESN is the most efficient.} Near-identical accuracy to TCN (RMSE
        0.0134) achieved in 33 seconds of training --- 290$\times$ faster. The
        reservoir's dynamical nature naturally echoes the Lorenz attractor geometry.

  \item \textbf{HPC and AI are complementary, not competing.} HPC handles accurate
        physics simulation to the checkpoint; AI handles rapid prediction beyond it.
        This framing correctly positions both components.

  \item \textbf{Chaotic systems impose fundamental limits.} All models diverge in
        long rollout due to the Lyapunov exponent. This is a physical constraint,
        not a modelling failure.
\end{enumerate}

\newpage 

\subsection{Future Work}

\begin{itemize}
  \item \textbf{Physics-Informed Neural Networks (PINNs):} Embed Lorenz equations
        in the loss function to enforce physical consistency.
  \item \textbf{Multi-GPU scaling:} Extend CUDA with MPI for $10^6$-trajectory
        ensembles.
  \item \textbf{Transformer surrogates:} Recent work shows $<$2\% relative L2 error
        on Lorenz with easy-attention Transformers~\cite{raissi2019}.
  \item \textbf{Adaptive checkpoint:} Dynamically set the handoff point based on
        estimated local Lyapunov exponent.
  \item \textbf{Uncertainty quantification:} Ensemble ESNs for prediction intervals.
\end{itemize}

\bibliographystyle{plainnat}
\begin{thebibliography}{9}

\bibitem{lorenz1963}
Lorenz, E.~N. (1963).
\newblock Deterministic nonperiodic flow.
\newblock \textit{Journal of Atmospheric Sciences}, 20(2), 130--141.

\bibitem{hairer1993}
Hairer, E., N{\o}rsett, S.~P., \& Wanner, G. (1993).
\newblock \textit{Solving Ordinary Differential Equations I} (2nd ed.).
\newblock Springer-Verlag.

\bibitem{raissi2019}
Raissi, M., Perdikaris, P., \& Karniadakis, G.~E. (2019).
\newblock Physics-informed neural networks.
\newblock \textit{Journal of Computational Physics}, 378, 686--707.

\bibitem{chen2018}
Chen, R.~T.~Q., Rubanova, Y., Bettencourt, J., \& Duvenaud, D. (2018).
\newblock Neural ordinary differential equations.
\newblock \textit{NeurIPS}, 31.

\bibitem{pathak2018}
Pathak, J., Hunt, B., Girvan, M., Lu, Z., \& Ott, E. (2018).
\newblock Model-free prediction of large spatiotemporally chaotic systems from
  data: A reservoir computing approach.
\newblock \textit{Physical Review Letters}, 120(2), 024102.

\bibitem{tcn2021}
Iwana, B.~K., \& Uchida, S. (2021).
\newblock An empirical evaluation of temporal convolutional networks for
  time series classification.
\newblock \textit{PLOS ONE}, 16(6), e0254841.

\bibitem{openmp}
OpenMP Architecture Review Board (2021).
\newblock \textit{OpenMP Application Programming Interface, Version 5.2}.
\newblock \url{https://www.openmp.org/}

\bibitem{nvidia_cuda}
NVIDIA Corporation (2023).
\newblock \textit{CUDA C++ Programming Guide}.
\newblock \url{https://docs.nvidia.com/cuda/cuda-c-programming-guide/}

\end{thebibliography}

% =============================================================
\appendix
\section{Complete Results Summary}
% =============================================================

\begin{center}
\rowcolors{2}{lightblue}{white}
\begin{tabular}{lrrrrrr}
  \toprule
  \rowcolor{accent!15}
  \textbf{Model} & \textbf{RMSE} & \textbf{MAE} &
  \textbf{Roll. Mean} & \textbf{Roll. Max} &
  \textbf{Train Time} & \textbf{Improvement} \\
  \midrule
  MLP (original) & 0.5256 & 0.4574 & 31.12 & 45.23 & --- & baseline \\
  MLP (improved) & 0.4027 & 0.3184 & 9.54  & 28.16 & 33 min  & +23.4\% \\
  LSTM (improved)& 0.2179 & 0.1643 & 8.54  & 22.18 & 41 min  & +58.5\% \\
  TCN            & \good{0.0128} & \good{0.0079} & \good{2.61} & \good{13.74} & 159 min & \good{+97.6\%} \\
  ESN            & 0.0134 & 0.0096 & 5.80  & 26.07 & \good{33 sec} & +97.5\% \\
  \bottomrule
\end{tabular}
\end{center}

\section{Master Notebook Structure}

Both notebooks are fully reproducible:
\begin{itemize}
  \item \textbf{lorenz\_hpc.ipynb} (27 cells) --- local HPC benchmarking,
        original MLP/LSTM, report generation.
  \item \textbf{lorenz\_improved.ipynb} (22 cells) --- Kaggle T4 GPU,
        improved LSTM + TCN + ESN, all plots.
\end{itemize}

\end{document}